无量纲聚焦立方非线性Schrödinger方程
\(i\psi_t+\frac12\psi_{xx}+|\psi|^2\psi=0\) 定义在全实轴，孤子解在无穷远消失。定义粒子数 \(N=\int|\psi|^2dx\)、动量 \(P=\int\operatorname{Im}(\psi^*\psi_x)dx\) 和Hamilton量 \(H=\int(\frac12|\psi_x|^2-\frac12|\psi|^4)dx\)。
(1) 使用驻波假设求精确亮孤子场，校验边界条件及一个积分形式的一阶运动方程。
(2) 独立计算孤子的 \(N,P,H\) 和能量标度；施以Galilean变换并核对移动孤子中附加动能。
(3) 对匀幅背景 \(\psi=Ae^{iA^2t}\) 做复扰动线性化，推导色散关系及模不稳定的临界波数和最快增长率。
(4) 对比孤子局域解和不稳定均匀波背景的守恒条件，讨论一维中的尺度平衡、\(\eta\to0\) 极限及为何不能把模不稳定阈值理解成孤子的自身失稳。
